% =====================================================================
%  Основное содержимое отчёта.
% =====================================================================

\labpart{Описание алгоритмов}

\subsection*{Свой алгоритм (лабораторная работа №3.1)}

Раздел «In-memory data» приложения \texttt{SecureApp}
(\texttt{Lab2.1-SecureApp}, добавлен в лабораторной работе №3.1):
\texttt{InMemoryRecordStore} на \texttt{ConcurrentDictionary} с выдачей
идентификаторов через \texttt{Interlocked.Increment}, конфиденциальные
поля шифруются AES-256-GCM (\texttt{ConfidentialCrypto}) непосредственно
при записи, вся валидация входных данных выполняется декларативно
средствами ASP.NET Core (\texttt{DataAnnotations} на DTO,
автоматическая обработка ошибок разбора JSON). Подробное описание и
обоснование см. в отчёте по лабораторной работе №3.1.

\subsection*{Предоставленный алгоритм}

\texttt{Lab3.2-ProvidedAlgorithm} -- отдельная консольная программа на
том же языке (C\#), решающая ту же задачу без использования ASP.NET
Core: сырой \texttt{System.Net.HttpListener}, обычный (не потокобезопасный)
\texttt{Dictionary}, ручной разбор JSON без описания схемы и без
валидации, «шифрование» через XOR с ключом, зашитым в исходный код
(\autoref{lst:provided-xor}), и постоянно растущий список
\texttt{activityLog}, в который записывается открытый текст каждой
операции (\autoref{lst:provided-log}) -- предположительно для отладки, но
без учёта того, что содержимое этого списка тоже является
конфиденциальными данными.

\begin{codelisting}[fontsize=\footnotesize]{csharp}{«Шифрование» через XOR с ключом, зашитым в исходный код}{provided-xor}
string XorObfuscate(string text, bool alreadyBase64 = false)
{
    var key = "my-static-secret-key-123"u8.ToArray();
    var bytes = alreadyBase64 ? Convert.FromBase64String(text) : Encoding.UTF8.GetBytes(text);
    for (var i = 0; i < bytes.Length; i++) bytes[i] ^= key[i % key.Length];
    return alreadyBase64 ? Encoding.UTF8.GetString(bytes) : Convert.ToBase64String(bytes);
}
\end{codelisting}

\begin{codelisting}[fontsize=\footnotesize]{csharp}{Открытый текст конфиденциального значения, записываемый в постоянно растущий журнал}{provided-log}
activityLog.Add($"{DateTime.UtcNow:o} CREATE confidential: title='{title}' content='{content}'");

var id = nextConfId++;               // не атомарно
confidential[id] = cipher;           // обычный Dictionary, без блокировки
\end{codelisting}

Единственная обработка ошибок во всей программе -- один блок
\texttt{catch (Exception ex)} на верхнем уровне обработчика запроса,
возвращающий текст исключения клиенту с кодом \texttt{500} независимо от
причины сбоя.

\labpartbreak{Анализ безопасности предоставленного алгоритма}

\subsection*{1. Обработка невалидных данных}

Тот же набор из семи проверок, что применялся к своему алгоритму в
лабораторной работе №3.1, был подан на предоставленный алгоритм
(\autoref{lst:fault-comparison}). Результат разителен: там, где свой
алгоритм для каждой из семи причин возвращает точный, предсказуемый код
ответа (\texttt{400} с описанием конкретной причины или корректный
\texttt{201}), предоставленный алгоритм отвечает \texttt{500} с сырым
текстом исключения на четыре из семи проверок -- вызывающая сторона не
может отличить «неверный тип поля» от «пустое тело запроса» иначе как
разбором текста исключения, а само это сообщение дополнительно
раскрывает внутреннее устройство сервера (пространства имён, сигнатуры
методов -- классический пример избыточно подробного сообщения об
ошибке). Ограничения на размер входных данных нет вовсе: содержимое из
10\,050 символов и содержимое размером 5~МБ приняты одинаково успешно
(\texttt{201}), что делает сервер уязвимым к исчерпанию памяти при
достаточном числе таких запросов.

\codelistingfile[fontsize=\scriptsize]{text}{lab/code/fault_injection_comparison.txt}{Обработка некорректного ввода: свой алгоритм против предоставленного}{fault-comparison}

\subsection*{2. Потокобезопасность}

200 одновременных запросов на создание записи (по 50 в одновременной
обработке) были поданы на оба алгоритма
(\autoref{lst:concurrency-comparison}). Свой алгоритм обработал все 200
без потерь: 200 уникальных идентификаторов, 200 сохранённых записей.
Предоставленный алгоритм вернул 200 ответов, но лишь 199 уникальных
идентификаторов -- два разных запроса получили один и тот же
идентификатор из-за неатомарного \texttt{nextConfId++}, и лишь 198
записей реально сохранились: одна запись была молча перезаписана другой
с тем же идентификатором. Это не сбой с сообщением об ошибке, а тихая
потеря данных, обнаруживаемая только сравнением числа отправленных и
реально сохранённых записей.

\codelistingfile[fontsize=\footnotesize]{text}{lab/code/concurrency_comparison.txt}{Состояние гонки при 200 параллельных запросах}{concurrency-comparison}

\subsection*{3. Криптографическая стойкость «шифрования»}

XOR с фиксированным, зашитым в исходный код ключом обратим тривиально:
доказательство -- независимая реализация той же операции на Python
(\autoref{lst:xor-poc}), использующая только ключ, скопированный из
исходного текста программы, без обращения к запущенному серверу или его
памяти. Такое «шифрование» не даёт защиты, если исходный код или
скомпилированный файл когда-либо станут доступны стороннему лицу --
ровно тот сценарий, для защиты от которого шифрование конфиденциальных
данных вообще применяется. Дополнительно, в отличие от AES-GCM,
используемого в своём алгоритме, у этой схемы нет ни одноразового
вектора (nonce) -- одинаковый открытый текст всегда даёт одинаковый
шифротекст, -- ни проверки подлинности: подмена байт шифротекста
напрямую в структуре данных не будет обнаружена при расшифровке.

\begin{codelisting}[fontsize=\footnotesize]{python}{Независимое воспроизведение и обращение XOR-«шифрования» по одному лишь ключу из исходного кода}{xor-poc}
KEY = b"my-static-secret-key-123"  # copied verbatim from Program.cs

def xor(data, key):
    return bytes(b ^ key[i % len(key)] for i, b in enumerate(data))
\end{codelisting}

\codelistingfile{text}{lab/code/xor_reversal_output.txt}{Результат восстановления открытого текста без обращения к серверу}{xor-poc-output}

\subsection*{4. Работа с конфиденциальными данными в оперативной памяти}

Дампы обоих алгоритмов снимались по методике, аналогичной лабораторной
работе №3.1, но с уточнением: поиск утилитой \texttt{strings} велся в
двух режимах -- 8-битном (находит только байтовые, как правило
UTF-8, буферы) и \texttt{-e~l} (UTF-16LE, необходим, чтобы увидеть сами
объекты \texttt{System.String} в управляемой куче .NET, которые в
8-битном режиме невидимы). Это уточнение самой методики поиска -- не
использованное в лабораторной работе №3.1 -- позволило впервые различить
временные буферы фреймворка от долгоживущих объектов приложения
(\autoref{lst:memory-comparison}).

\codelistingfile[fontsize=\scriptsize]{text}{lab/code/memory_dump_comparison.txt}{Число вхождений маркера до/после удаления и принудительной сборки мусора}{memory-comparison}

У своего алгоритма число вхождений маркера не изменилось после удаления
записи и принудительной полной сборки мусора (3 и 2 -- в 8-битном и
16-битном режиме соответственно, для обоих дампов одинаково) -- все
найденные копии оказались буферами конвейера ASP.NET Core, а не
структурами данных самого хранилища. У предоставленного алгоритма число
вхождений заметно уменьшилось после сборки мусора (с 6 до 2 и с 7 до 2
соответственно) -- то есть большая часть временных копий действительно
была собрана, -- но два вхождения остались навсегда: это оказались
буквальные строки из \texttt{activityLog}, что было подтверждено прямым
чтением найденных строк из дампа. Иначе говоря, у предоставленного
алгоритма после удаления записи и явной сборки мусора конфиденциальное
значение остаётся в памяти процесса из-за собственной, легко устранимой
ошибки кода (логирование секрета), тогда как у своего алгоритма остаток
объясняется исключительно платформой и не устраним на уровне кода
приложения без отказа от использующегося фреймворка целиком.

\labpartbreak{Разработанная методика оценки безопасности кода}

Методика построена как контрольный список из шести критериев, каждый --
с конкретной, воспроизводимой техникой проверки, выведенной из
практического опыта лабораторных работ №2.1--3.1, и трёхуровневой
оценкой (соответствует / частично соответствует / не соответствует).

\begin{enumerate}
  \item \textbf{Устойчивость к невалидному вводу.} Подать не менее шести
  видов заведомо некорректного ввода (неверный тип, синтаксическая
  ошибка, отсутствующее поле, пустое тело, превышение разумного размера,
  управляющие символы) и проверить: (а) отклоняется ли каждый некорректный
  случай явным, специфичным для причины кодом ответа, а не общим
  сбоем; (б) не раскрывается ли во внешнем ответе внутренняя
  информация (текст исключения, путь к файлу, версия фреймворка).
  \item \textbf{Потокобезопасность разделяемого состояния.} Подать не
  менее 100--200 одновременных операций записи и сравнить число
  отправленных запросов, число уникальных возвращённых идентификаторов и
  число реально сохранившихся записей -- любое расхождение означает
  состояние гонки.
  \item \textbf{Криптографическая стойкость защиты конфиденциальных
  данных.} Определить алгоритм и источник ключа по коду; проверить
  наличие одноразового вектора/соли на каждую операцию (то есть
  разные ли шифротексты получаются из одинакового открытого текста) и
  наличие проверки подлинности (защиты от подмены шифротекста);
  попытаться независимо, без обращения к работающему серверу,
  воспроизвести обратное преобразование по одному лишь исходному коду.
  \item \textbf{Устойчивость к нагрузке по размеру входных данных.}
  Определить, существует ли явный верхний предел размера принимаемых
  данных, и подать вход, заметно превышающий разумный размер (на
  два--три порядка), проверив код ответа и то, что процесс не потребляет
  ресурсы неограниченно.
  \item \textbf{Работа с конфиденциальными данными в оперативной
  памяти.} Снять дамп процесса после создания, после изменения и после
  удаления конфиденциальной записи с уникальным маркером; выполнить
  двухрежимный поиск маркера в каждом дампе (8-битный и
  \texttt{strings~-e~l}, то есть UTF-16LE) до и после принудительной
  полной сборки мусора. Если после сборки мусора маркер остаётся --
  установить, является ли источник кодом самого приложения
  (структуры данных, журналы) или инфраструктурой платформы, поскольку
  от этого зависит, устранима ли находка в принципе на уровне кода
  приложения.
  \item \textbf{Изоляция сбоев.} Убедиться, что некорректный запрос не
  влияет на обработку последующих корректных запросов и не приводит к
  зависанию или падению процесса в целом.
\end{enumerate}

Обоснование корректности методики: каждый критерий сформулирован как
воспроизводимая, автоматизируемая проверка с однозначно интерпретируемым
результатом (числовое сравнение или бинарный факт «найдено/не найдено»),
а не как экспертная оценка «на глаз» -- поэтому применение методики к
одному и тому же алгоритму разными людьми в разное время должно давать
одинаковый результат. Дополнительно методика была опробована на двух
существенно разных по качеству реализациях, продемонстрировав
уверенное разделение (см. далее) -- то есть не является тривиальной или
всегда дающей одинаковый ответ вне зависимости от входа.

\labpartbreak{Сравнение своего варианта с предоставленным по методике}

\begin{table}[H]
\centering
\small
\begin{tabular}{|p{4.3cm}|p{4.0cm}|p{4.6cm}|}
\hline
\textbf{Критерий} & \textbf{Свой алгоритм (№3.1)} & \textbf{Предоставленный алгоритм} \\
\hline
1. Невалидный ввод & Соответствует (7/7 точных кодов, 0 утечек деталей) & Не соответствует (4/7 -- общий 500 с текстом исключения) \\
\hline
2. Потокобезопасность & Соответствует (0 потерь из 200) & Не соответствует (2 записи из 200 потеряны молча) \\
\hline
3. Криптостойкость & Соответствует (AES-256-GCM, nonce, тег подлинности) & Не соответствует (обратимый XOR, статический ключ, без nonce и без тега) \\
\hline
4. Пределы размера входа & Соответствует (400 уже на 10\,050 симв.) & Не соответствует (принят вход 5~МБ без ограничений) \\
\hline
5. Конфиденциальность в памяти & Частично (остаток есть, но не устраним и не является ошибкой кода) & Не соответствует (остаток -- прямое следствие собственной ошибки: логирование секрета) \\
\hline
6. Изоляция сбоев & Соответствует (сервер работоспособен после всех проверок) & Соответствует (сервер работоспособен после всех проверок, хотя и ценой недифференцированных 500) \\
\hline
\end{tabular}
\caption{Сравнение по разработанной методике}
\end{table}

Из шести критериев предоставленный алгоритм полностью соответствует
только одному (изоляция сбоев -- сервер не падает целиком, хотя и
обрабатывает ошибки крайне грубо), частично или полностью не
соответствует остальным пяти. Свой алгоритм полностью соответствует
пяти критериям из шести и лишь частично -- шестому, причём именно
потому, что этот шестой пункт (остаточное присутствие данных в памяти
после сборки мусора) для управляемой платформы в принципе неустраним
средствами прикладного кода, что и зафиксировано методикой отдельно от
устранимых нарушений.
